\begin{recipe}{Puffy pita}
\kicker{Brood · zelf bakken}
\index[register]{Puffy pita}

\meta{6 stuks \dvd Rijzen 1 uur \dvd Oven 240\,°C}

\lettrine{Z}{elfgebakken} pita die in een paar minuten opbolt tot een luchtig,
hol broodje. Plakkerig deeg, een gloeiend hete steen — en het komt helemaal
goed. Perfect om in de shakshuka te dippen.

\blockrule{Ingrediënten}
\begin{ingredients}
  \ing{300 g}{bloem}
  \ing{7 g}{instant gist}
  \ing{200 ml}{lauw water}
  \ing{1 el}{suiker}
  \ing{1 el}{olijfolie}
  \ing{2 tl}{zout}
\end{ingredients}

\blockrule{Bereiding}
\tip{Deeg te plakkerig? Niet panikeren — een extra lepel bloem en doorkneden,
het komt goed.}
\begin{steps}
  \item Roer gist en suiker door het lauwe water en laat het 8–9 minuten tot
        leven komen.
  \item Meng de bloem met het zout in een grote kom.
  \item Roer de olijfolie door het gistmengsel, giet alles bij de bloem en meng
        tot er een deeg ontstaat.
  \item Kiep het deeg op de werkbank en kneed 2–3 minuten. Ja, het is plakkerig —
        en ja, het komt goed; gooi voor de zekerheid nog 1 el bloem erbij. Doe
        het in een ingevette kom, dek af met (ingeolied) folie en laat 1 uur
        rijzen op een warme plek.
  \item Zet een bakplaat of pizzasteen in de oven en verhit minstens 1 uur op
        230–240\,°C. Hoe heter de steen, hoe mooier je pita's straks opbollen.
  \item Sla de lucht uit het deeg, verdeel in 6 gelijke stukken en bol op. Leg op
        een bebloemde plaat en laat een half uur rusten.
  \item Rol elk bolletje uit tot een schijf van 15 cm en leg op een passend
        stukje bakpapier.
  \item Bak de pita's 2–3 minuten op de hete plaat tot ze mooi puffy worden.
\end{steps}
\end{recipe}
